\subsection{Neural Language Models}
\label{sec:neural-language-models}
\slidesrc{Language models \& RNNs: 35--50}

n-gram models fail because they treat every word sequence as a separate, unrelated
item: seeing ``she went to school'' a thousand times teaches them nothing about ``he
walked to college''. Neural language models fix this by representing words as
\textbf{vectors}, so that similar words get similar representations and what is learnt
about one word transfers to its neighbours.

\subsubsection{Aside: word embeddings via word2vec}
\slidesrc{Language models \& RNNs: 36}

\paragraph{One-hot vectors.}
\notinslides
A computer needs numbers, not words. The simplest way to turn a word into numbers is a
\textbf{one-hot vector}: number the words of the vocabulary from 1 to $|V|$; word number
$i$ becomes a vector of length $|V|$ with a 1 in position $i$ and 0 everywhere else.
With the vocabulary $(\textit{cat}, \textit{dog}, \textit{the}, \textit{sleeps})$:
\[
  \textit{cat} = \col{1\\0\\0\\0}, \quad \textit{dog} = \col{0\\1\\0\\0}, \quad
  \textit{the} = \col{0\\0\\1\\0}, \quad \textit{sleeps} = \col{0\\0\\0\\1}.
\]
One-hot vectors identify words but say nothing about their meaning: \emph{cat} is
exactly as different from \emph{dog} as from \emph{sleeps} (any two of them have
dot product 0).

\paragraph{Word embeddings.}
The goal of a \textbf{word embedding} is to represent each word with a short, dense
vector (a few hundred real numbers) such that:
\begin{itemize}[nosep]
  \item the vector of a word is \textbf{similar} to the vectors of words that occur in
        \emph{similar} contexts;
  \item it is \textbf{dissimilar} from the vectors of words that are unlikely to appear
        in the same contexts.
\end{itemize}
This is the distributional hypothesis (Section~\ref{sec:distributional-hypothesis})
turned into vectors (Section~\ref{sec:geometric-interpretation}). \textbf{Word2vec} is
a popular method that learns such vectors by training a small neural network to
predict words from their neighbours. Similarity between two vectors is measured by
their \textbf{dot product}: multiply the entries position by position and add.

\begin{nbexample}{similar and dissimilar words}
\slidesrc{Language models \& RNNs: 36}
Eight-dimensional embeddings from the slide:
\[
  \textit{baking} = \col{0.286\\0.792\\-0.177\\-0.107\\0.109\\-0.542\\0.349\\0.271}\!,\quad
  \textit{flour}  = \col{0.413\\0.582\\-0.007\\0.247\\0.216\\-0.718\\0.147\\0.051}\!,\quad
  \textit{dollar} = \col{-0.154\\0.024\\0.512\\0.782\\0.359\\0.141\\-0.581\\-0.027}\!.
\]
\begin{align*}
  \textit{baking}\T\textit{flour}
  &= 0.286 \cdot 0.413 + 0.792 \cdot 0.582 + (-0.177)(-0.007) + \cdots\\
  &\qquad + 0.271 \cdot 0.051\\
  &= 0.118 + 0.461 + 0.001 - 0.026 + 0.024 + 0.389 + 0.051 + 0.014\\
  &\approx 1.03
\end{align*}
The large positive entries are in the same places (positions 2 and 6 contribute most),
so the dot product is large: \emph{baking} and \emph{flour} appear in similar contexts
(recipes). For \emph{baking} and \emph{dollar} the signs disagree in most positions and
the dot product is negative: they are used in unrelated contexts.
\end{nbexample}

\begin{nbwarning}[the number on the slide]
The slide gives $\textit{baking}\T\textit{dollar} = -0.35$. With the vectors as printed
the dot product is $-0.45$; $-0.34$ is their \emph{cosine similarity} (the dot product
divided by the lengths of the two vectors, Section~\ref{sec:geometric-interpretation}).
Either way the conclusion is the same: negative, hence dissimilar.
\end{nbwarning}

\subsubsection{A fixed-window neural language model}
\slidesrc{Language models \& RNNs: 37--41}

Recall the language-modelling task: given a sequence of words
$\vect x^{(1)}, \ldots, \vect x^{(t)}$, output the probability distribution of the next
word, $p(\vect x^{(t+1)} \mid \vect x^{(1)}, \ldots, \vect x^{(t)})$. The first neural
approach keeps the n-gram idea of a \textbf{fixed window} of previous words, but
replaces counting with a \textbf{neural network} that \emph{learns} to predict.

For ``Anna was tired so she went to her \underline{\quad}'', a window of 4 words
takes \emph{she went to her} and predicts the missing word. The network is built in
four layers, from bottom to top:

\begin{center}
\begin{tikzpicture}[>=Stealth,
    neu/.style={circle, draw=#1, fill=#1!35, minimum size=5pt, inner sep=0pt},
    lay/.style={draw=#1, rounded corners, thick, inner sep=2pt},
    lab/.style={font=\small, align=left, anchor=west, text width=4.3cm}]
  % one-hot words
  \foreach \w [count=\i] in {she, went, to, her} {
    \node[font=\small] (w\i) at (1.6*\i,0) {\w};
    \node[font=\scriptsize, nbgray] at (1.6*\i,-0.4) {$\vect x^{(\i)}$};
    % embedding block
    \foreach \k in {1,...,4} \node[neu=magenta] (e\i\k) at (1.6*\i+0.24*\k-0.6,1.2) {};
    \draw[->, nbgray] (w\i) -- node[right, font=\scriptsize] {$E$} (1.6*\i,1.05);
  }
  \node[lay=magenta, fit=(e11)(e44)] (emb) {};
  % hidden
  \foreach \k in {1,...,10} \node[neu=cyan] (h\k) at (2.9+0.32*\k,2.5) {};
  \node[lay=cyan, fit=(h1)(h10)] (hid) {};
  \draw[->, thick] (emb) -- node[right, font=\scriptsize] {$W_e$} (hid);
  % output histogram
  \node[lay=nbred, minimum width=5.6cm, minimum height=0.45cm, font=\scriptsize]
       (out) at (4.65,3.8) {distribution over the whole vocabulary};
  \draw[->, thick] (hid) -- node[right, font=\scriptsize] {$U$} (out);
  \foreach \x/\h in {0.2/0.15,0.5/0.4,0.8/0.2,1.3/0.6,1.6/0.25,2.1/0.3,2.6/0.1,
                     3.0/0.8,3.3/0.35,3.9/0.2,4.4/0.12,4.9/0.3}
    \fill[nbred!70] (1.85+\x,4.05) rectangle ++(0.15,\h);
  % labels
  \node[lab, nbred]   at (8.0,3.85) {\textbf{output}\\$\hat{\vect y} = \softmax(U\vect h) \in \mathbb R^{|V|}$};
  \node[lab, cyan!60!black] at (8.0,2.5) {\textbf{hidden layer}\\$\vect h = \sigma(W_e \vect e)$};
  \node[lab, magenta] at (8.0,1.2) {\textbf{word embeddings}\\$\vect e = [\vect e^{(1)}, \vect e^{(2)}, \vect e^{(3)}, \vect e^{(4)}]$};
  \node[lab]          at (8.0,-0.15) {\textbf{one-hot word vectors}\\$\vect x^{(1)}, \vect x^{(2)}, \vect x^{(3)}, \vect x^{(4)}$};
\end{tikzpicture}
\captionof{figure}{The fixed-window neural language model. Information flows upwards:
four words $\to$ their embeddings, glued together $\to$ a hidden layer $\to$ a
probability for every word of the vocabulary (tall bars = likely next words).}
\end{center}

\begin{enumerate}
  \item \textbf{One-hot word vectors} $\vect x^{(1)}, \ldots, \vect x^{(4)} \in
        \mathbb R^{|V|}$: the four words of the window.
  \item \textbf{Word embeddings.} Each word is turned into its embedding
        $\vect e^{(i)} = E \vect x^{(i)}$. $E$ is a matrix with one column per word of
        the vocabulary ($E \in \mathbb R^{d \times |V|}$, $d$ = embedding size).
        Multiplying a matrix by a one-hot vector simply \emph{picks one column}: the
        column of that word. So $E$ is just a table of word vectors, stored as a
        matrix. The four embeddings are then \textbf{concatenated} (written one after
        the other) into a single vector $\vect e$ of length $4d$.
  \item \textbf{Hidden layer} $\vect h = \sigma(W_e \vect e)$. The matrix $W_e$ mixes all
        the numbers of the four embeddings; $\sigma$ is a non-linear function applied
        to each entry (see below). $\vect h$ is a summary of the window useful for
        prediction.
  \item \textbf{Output} $\hat{\vect y} = \softmax(U \vect h) \in \mathbb R^{|V|}$: one
        number per word of the vocabulary, turned into a probability distribution by
        the softmax (see below). $\hat{\vect y}_w$ is the predicted probability that
        the next word is $w$. The hat in $\hat{\vect y}$ means ``predicted''.
\end{enumerate}

\notinslides
Two building blocks deserve a closer look.
\begin{itemize}
  \item \textbf{The non-linearity $\sigma$.} Typically the \emph{sigmoid}
        $\sigma(z) = \frac{1}{1 + e^{-z}}$, which squashes any number into $(0, 1)$, or
        $\tanh$, which squashes into $(-1, 1)$. Without it, the whole network would be
        a product of matrices, i.e.\ just one big linear map, and could only learn
        very simple relations. (Real networks also add a \emph{bias} vector inside
        $\sigma$, $\sigma(W_e \vect e + \vect b)$; the slides omit it for simplicity.)
  \item \textbf{The weights $E$, $W_e$, $U$} are the \textbf{parameters} of the model,
        written together as $\theta$. They start as random numbers; \emph{training}
        (below) adjusts them so that the predictions become good.
\end{itemize}

\subsubsection{How to train a neural language model}
\slidesrc{Language models \& RNNs: 42--49}

\begin{nbdef}{The four ingredients of training a neural network}
\begin{enumerate}[nosep]
  \item \textbf{A corpus of data}: here a collection of documents, i.e.\ a long sequence
        of words $\vect x^{(1)}, \ldots, \vect x^{(T)}$.
  \item \textbf{A model}: here the fixed-window neural LM.
  \item \textbf{A loss function $J$}: a number measuring how wrong the model's
        predictions are; it compares the output distribution with the true one.
  \item \textbf{An optimisation algorithm} that changes the parameters to make the
        loss smaller: \textbf{stochastic gradient descent}.
\end{enumerate}
\end{nbdef}

\paragraph{The softmax.}
The vector $U\vect h \in \mathbb R^{|V|}$ contains one \emph{score} per word, but scores
can be negative and do not add up to 1: it is an \textbf{un-normalised} probability
distribution. The \textbf{softmax} normalises it:
\[
  \softmax(\vect z)_i = \frac{e^{z_i}}{\sum_{j=1}^{K} e^{z_j}}, \qquad i = 1, \ldots, K .
\]
The exponential makes every entry positive (and enlarges the gaps between scores),
and dividing by the sum makes the entries add up to 1.

\begin{nbexample}{softmax}
\slidesrc{Language models \& RNNs: 47}
Scores $U\vect h = (1.3,\ 5.1,\ 2.2,\ 0.7,\ 1.1)$ for a vocabulary of $K = 5$ words.
\begin{align*}
  (e^{1.3}, e^{5.1}, e^{2.2}, e^{0.7}, e^{1.1})
  &= (3.67,\ 164.02,\ 9.03,\ 2.01,\ 3.00) \reason{exponentiate}\\
  \text{sum} &= 3.67 + 164.02 + 9.03 + 2.01 + 3.00 = 181.73 \\
  \hat{\vect y} &= (0.02,\ 0.90,\ 0.05,\ 0.01,\ 0.02) \reason{divide by the sum}
\end{align*}
Check: $0.02 + 0.90 + 0.05 + 0.01 + 0.02 = 1$. The highest score (5.1) gets most of the
probability, but no word gets exactly 0.
\end{nbexample}

\paragraph{The loss: cross-entropy.}
We compare
\begin{itemize}[nosep]
  \item the \textbf{computed} probability $\hat y_i = q(V_i \mid \textit{she went to
        her})$, what the model says about word $V_i$ of the vocabulary;
  \item the \textbf{true} probability $y^*_i = p(V_i \mid \text{full context})$, what
        the word really is.
\end{itemize}
The standard way to compare two distributions is the \textbf{cross-entropy}:
\[
  \mathrm{CE}(\vect y^*, \hat{\vect y}) = - \sum_{i=1}^{|V|} y^*_i \log \hat y_i .
\]
In training data we do not know the ``true distribution'': we only know which word
\emph{actually came next} in the corpus. So $\vect y^*$ is a \textbf{one-hot vector}:
1 for the true next word, 0 for every other word. All the terms of the sum vanish
except one. If the true word in the data is the fifth word of the sentence,
$\vect x^{(5)}$, then
\[
  J^{(t)} = \mathrm{CE}(\vect y^*, \hat{\vect y}) = -\log \hat y_{\vect x^{(5)}}
  \qquad \text{(the negative log-probability of the true word).}
\]
(The slide writes $\hat y_5$, meaning ``the entry of $\hat{\vect y}$ belonging to the
true word $\vect x^{(5)}$''.)

\begin{nbintuition}[why $-\log$ of the probability?]
\notinslides
$-\log p$ is 0 when $p = 1$ (the model was certain and right: no penalty) and grows
without limit as $p \to 0$ (the model said the true word was nearly impossible: huge
penalty). With the softmax output above (natural logarithm):
\begin{itemize}[nosep]
  \item if the true word is the 2nd one ($\hat y = 0.90$): loss
        $= -\ln 0.90 \approx 0.10$, small;
  \item if the true word is the 4th one ($\hat y = 0.011$): loss
        $= -\ln 0.011 \approx 4.50$, large.
\end{itemize}
Minimising the loss means: \emph{give high probability to the words that really come
next}.
\end{nbintuition}

\paragraph{Gradient descent.}
\notinslides
The loss $J(\theta)$ depends on all the parameters $\theta$ (the entries of $E$, $W_e$,
$U$). The \textbf{gradient} $\nabla_\theta J$ is the vector of the derivatives of $J$
with respect to each parameter: it points in the direction in which the loss
\emph{increases} fastest. \textbf{Gradient descent} takes a small step in the opposite
direction, again and again:
\[
  \theta \leftarrow \theta - \alpha\, \nabla_\theta J(\theta),
\]
where $\alpha > 0$ (the \textbf{learning rate}) controls the step size. It is like
walking down a hill in the fog: you cannot see the valley, but you can feel the slope
under your feet and step downhill. The gradients are computed automatically by
\textbf{backpropagation}, i.e.\ the chain rule of calculus applied layer by layer from
the output back to the input. \emph{Stochastic} gradient descent computes the gradient
on a small random part of the data at each step instead of on the whole corpus
(Section~\ref{sec:recurrent-neural-networks}).

\subsubsection{Pros and cons of the fixed-window model}
\slidesrc{Language models \& RNNs: 50}

This model was introduced by Y.~Bengio et al., \emph{A Neural Probabilistic Language
Model}, 2003.

\textbf{Improvements over n-gram LMs:}
\begin{itemize}[nosep]
  \item \textbf{no sparsity problem}: the network computes a probability for every word
        from the embeddings, so unseen sequences get sensible, non-zero probabilities;
        if \emph{school} and \emph{college} have similar embeddings, what is learnt
        for one transfers to the other;
  \item \textbf{no need to store n-grams}: only the weights are stored.
\end{itemize}

\textbf{Remaining problems:}
\begin{itemize}[nosep]
  \item the \textbf{fixed window is too restrictive}: like an n-gram model it cannot see
        anything older than its window (``Anna was \emph{tired}'' is still invisible);
  \item \textbf{enlarging the window increases the size of $W_e$}: $W_e$ has one column
        per entry of $\vect e$, and $\vect e$ grows with every extra word; a window big
        enough for any text is impossible;
  \item \textbf{no symmetry}: $\vect x^{(1)}$ and $\vect x^{(2)}$ are multiplied by
        \emph{different} columns of $W_e$. What the model learns about the word
        \emph{school} in position 4 does not transfer to \emph{school} in position 2:
        the same knowledge has to be learnt separately for every position.
\end{itemize}

\begin{nbwhy}
The question left open by this slide, \emph{how can we process inputs of arbitrary
length?}, is answered by recurrent neural networks (next section), and later by
transformers. Both keep the two good ideas introduced here: words as learnt
embeddings, and a softmax over the vocabulary trained with cross-entropy. Every modern
LLM, GPT included, still ends with exactly this softmax layer.
\end{nbwhy}

\begin{nbsummary}[neural language models]
\begin{itemize}[nosep]
  \item Embedding: $\vect e = E\vect x$ (picks the column of word $\vect x$); similar
        words have a large dot product.
  \item Fixed window: $\hat{\vect y} = \softmax(U\, \sigma(W_e [\vect e^{(1)}, \ldots,
        \vect e^{(n)}]))$.
  \item $\softmax(\vect z)_i = e^{z_i} / \sum_j e^{z_j}$: positive, sums to 1.
  \item Loss: cross-entropy with a one-hot target $= -\log \hat y_{\text{true word}}$.
  \item Training: corpus + model + loss + (stochastic) gradient descent
        $\theta \leftarrow \theta - \alpha \nabla_\theta J$.
  \item Solves sparsity and storage; still a fixed window, growing $W_e$, no weight
        sharing across positions.
\end{itemize}
\end{nbsummary}
